\section{模型假设}
模型假设只用于补足题面未给出的物理细节，不能替代题面约束。表~\ref{tab:assumptions} 同时给出假设依据及其对模型的影响，便于在结果讨论中追踪适用范围。

\begin{table}[H]
\centering
\caption{模型假设、合理性依据及影响}
\label{tab:assumptions}
\begin{tabularx}{\textwidth}{C{0.8cm}L L L}
\toprule
编号 & 假设内容 & 合理性依据 & 对模型的影响 \\
\midrule
A1 & 题面给出的换能器开角、坡度、中心水深和附件深度数据真实有效。 & 题目未提供测量误差分布，直接使用给定数据是唯一可复现的输入口径。 & 模型输出为给定数据条件下的确定性结果，不包含观测误差置信区间。 \\
A2 & 问题一至问题三的海底为均匀平面坡面，声束扇面始终垂直于测线。 & 与题面理想海底描述一致，且多波束覆盖宽度由横航向剖面决定。 & 可将三维几何化为二维直线交点问题，并获得显式公式。 \\
A3 & 问题四在网格单元内使用双线性插值，局部梯度取该插值面的解析偏导。 & 原始网格步长为 $0.02$ 海里，双线性插值不额外制造高频地形。 & 网格内深度与梯度连续分片可计算，但可能平滑小于原网格尺度的起伏。 \\
A4 & 忽略声速剖面、船体姿态、安装误差、海况和障碍物。 & 题目没有提供相关参数或校准数据，任意赋值会降低结果可审计性。 & 得到的是几何测线方案，实际作业前仍需进行设备和海况修正。 \\
A5 & 问题四只要求测线本身位于区域内，测线之间的调头连接航段不计入题目规定的测线总长度。 & 题目指标针对测线总长度，未规定航线连接顺序和转场成本。 & 输出可直接回答题目指标，但不能替代完整航次调度方案。 \\
A6 & 题目未给出具体船型的最小转弯半径，曲线只要求切向连续且曲率有限。 & 无法从题面确定船舶动力学阈值。 & 报告最小曲率半径作为诊断值，不宣称满足某一船型。 \\
\bottomrule
\end{tabularx}
\end{table}
